%Metadata
\documentclass[12pt]{report}
\usepackage[english]{babel}
\usepackage[utf8x]{inputenc}
\usepackage{booktabs} %Table
\usepackage{multirow} %Rows
\usepackage{graphicx} %Insert Pictures
\usepackage{enumitem} %ToC
\usepackage{titlesec} 
\usepackage{geometry} %Margins

%=========== DISABLE THIS WHEN EXPORT... Remove the %
\setkeys{Gin}{draft=False}


\titleformat{\chapter}{\normalfont\huge}{\thechapter.}{20pt}{\huge\it}
\titlespacing*{\chapter}{0pt}{50pt}{40pt}
\titlespacing*{name=\chapter}{0pt}{-30pt}{10pt}

\renewcommand{\thetable}{\Alph{table}}
\title{Pro Mechanical Teletubbies}
\linespread{1.5}
\begin{document}


%----------------------------------------------------------------------------------------
%	TITLE PAGE
%----------------------------------------------------------------------------------------
	
\begin{titlepage}
\pagenumbering{roman} %Use roman letters for number at the start
\newcommand{\HRule}{\rule{\linewidth}{0.5mm}}

\center % Center everything on the page

%=================HEADING =======================================
\includegraphics{UoA_logo.jpg}\\
\textsc{\large School of Mechanical Engineering}\\[1cm]
\textsc{\Huge Design Practice}\\
\textsc{\large Mech Eng 2100}\\

%===============TITLE =======================

\HRule \\[0.4cm]
{ \Huge \bfseries Warman Competition}\\[0.2cm]
{ \Huge \bfseries Final Report}\\[0.1cm]
\HRule \\[-0.2cm]

%================NAMES ==========================

\huge TEAM 29 \\ [-0.2cm]
\Large The Professional Mechanical Teletubbies \\[1 cm]

\begin{minipage}{0.5\textwidth}
\begin{flushleft} \small
  \textsc{Xing Yong TAN}
   \linebreak
 \textsc{a1732497}
 \linebreak
  \linebreak
  \textsc{Shanmughanathan LAKSHMANAN}
   \linebreak 
 \textsc{a1741366}
  \linebreak
  \linebreak
  \textsc{Hamza Saquib Zain KHAN}
  \linebreak
  \textsc{a1741366}
 
\end{flushleft}
\end{minipage}
~
\begin{minipage}{0.4\textwidth}
\begin{flushright} \small
  \textsc{Peng Hoe HOR}
   \linebreak
 \textsc{a1732412}
 \linebreak
  \linebreak
  \textsc{Zheng Bing LIM}
   \linebreak
 \textsc{a1732649}
 \linebreak
 \linebreak
  \linebreak
\end{flushright}
\end{minipage}\\[0.5cm]
\end{titlepage}

%----------------------------------------------------------------------------------------
%	EXECUTIVE SUMMARY [done]
%----------------------------------------------------------------------------------------
\section*{Executive Summary}
This report details the design and build process of a robot for the Warman Design and Build Competition. The robot is required to move a payload from a pole to a bunker at the end of the competition track. Several designs for the movement of the robot is considered at part 2. They are the 8 wheel, 4 wheel and 2 Wheel design. For the robot to collect the payload on the top of the pole, a scissor lift design is considered. For the aforementioned designs, this report details the manufacturing process and testing processes. While the robot failed at its objective during its final run, we are confident that the robot will perform its objectives if more time for manufacturing and testing is given.
\setcounter{page}{2} %Set page number to 2 as title page dont count 
\pagebreak

%----------------------------------------------------------------------------------------
%	Table of Content [WIP]
%----------------------------------------------------------------------------------------
\section*{Table of Contents}
\begin{tabular}{ p{0.9\textwidth} p{0.1\textwidth} }

%TOC Details
\begin{enumerate}
    \item []Executive Summary 
    \item []Acknowledgements 
    \item Introduction 
	\begin{enumerate}
	\item [1.1] Problem specification
	\item [1.2]	Literature Review    
	\end{enumerate}
\item Design Alternatives 
	\begin{enumerate}
	\item [2.1]	Design A 
	\item [2.2]	Design B 
	\item [2.3]	Design C 
	\end{enumerate}
\item Final Design 
	\begin{enumerate}
	\item [3.1] Design Concept
    \item [3.2] Testing
    \item [3.3] Modifications
	\end{enumerate}
\item Conclusion
\item Reference List
\end{enumerate}

&
%Page numbers. Custom linewidth 
\begin{flushright}
i \\[0.2\linewidth]
ii \\[0.2\linewidth]
\end{flushright}
\end{tabular}
\pagebreak

\begin{tabular}{ p{0.9\textwidth} p{0.1\textwidth} }
Appendix
\begin{enumerate}[label=\Alph*.]
\item Project Planning and Management
\item Problem Solving and Team Reflection
\item Comparison Matrix 
\item Project Management Tools 
\item Supplementary Drawings 
\item Glossary of Terms 
\end{enumerate}
&
\begin{flushright}
\vspace{-0.3cm}
\item []  19 \\[0.25\linewidth]
\item []  20 \\[0.3\linewidth]
\end{flushright}

\tabularnewline
\end{tabular}
\pagebreak

%----------------------------------------------------------------------------------------
%	ACKNOWLEDGMENTS [Done]
%----------------------------------------------------------------------------------------
\section*{Acknowledgements}
We would like to thank our course coordinator, Dr. Antoni Blazewicz for giving us this wonderful opportunity to participate in this project giving us an experience in the field of robotics. Furthermore, we would like to thank head tutor Nathaniel Shearer and his team Saxon A. Nelson-Milton, Harrison Schuster, Thomas Jacquier, Joshua Leyshon for providing valuable guidance in designing and building our robot.

Also, many thanks go to Mark and Nathan for helping us manufacture our robot in Holden Labs. Last but not least, we would like to thank the organisers Engineers Australia, and sponsors Weir Minerals for organising the Warman Design and Build Competition.


\pagebreak

%----------------------------------------------------------------------------------------
%Introduction [WIP]
%----------------------------------------------------------------------------------------
\pagenumbering{arabic} %Change to numerals
\chapter{Introduction}

The objective of the robot is to transport the payload safely from a pole to the bunker at the end of the track. During travel, the payload cannot be dropped and has to maintain the correct orientation until it is dropped into the bunker.

The problem can be divided into 3 parts. The first of which is the robot needs to maneuver and travel to the correct position in the correct orientation on the competition track. The second part is the collecting mechanism of the payload There are 3 options/locations that the power packs are located all of which are of different heights and provide different points with the tallest providing the most, the medium height pole providing 2/3 and the shortest providing 1/3 of points. For first two poles the robot is required to extend up to specific heights to collect the payload, whereas for the shortest pole a surface level design is required to collect the payload. The final part of the problem is the release mechanism. The release mechanism must be able to deposit the payload with the arrow facing the correct direction to maximize points. There are steps before the bunker therefore the releasing mechanism must be able to overcome these steps to place the payload into the bunker. The method of placing the power pack itself needs to be considered as the bunker has a size of 150 x 150 mm, the power packs needs to be placed within these constraints to maximise the points per Warman competition rules. Placing it anywhere else will provide lesser points.


The solutions to the problem are further limited by the maximum weight (6kg) and maximum dimensions (400mmx400mm) per the competition rules. Moreover, the robot needs to designed within a tight budget of A\$50.

\pagebreak
 

\section{Literature review} 
 Wheels play an important part in a robot as it is needed for the motion of the robot and stability. However, there are many constraints and specifications needed to be considered while manufacturing a wheel system. For example, the motors required to drive wheels, the type of materials needed to manufacture the wheels and also other features that are needed to be added to maintain the balance of the whole robot.
 
First and foremost, a moment force needs to be applied to the wheels for motion. Hence, motors such as Direct Current (DC) motors, servos motors and stepper motors have to be considered to produce the similar type of force (Warren, Adams \& Molle 2011). There are generally two types of DC motors which are brushed DC motors and brushless DC motors (Warren, Adams \& Molle 2011). Brushed DC motors are easy to use as it needs to be physically connected to a power supply through both positive and ground connection (Warren, Adams \& Molle 2011). However, brushed DC motors consumes a high amount of current for the motors to work (Warren, Adams \& Molle 2011). If the workload increases, the DC motors will require more current for a higher speed (Warren, Adams \& Molle 2011). The brushless DC motors have similar properties but these motors have a longer life cycle and higher reliability because it uses a three-phase driver circuit instead (Warren, Adams \& Molle 2011). However, stepper motors have higher torque-to-weight ratio than DC motors (Greenough \& Kung 2013). Furthermore, stepper motors are more accurate in terms of rotations and position compared to DC motors (Greenough \& Kung 2013). To drive the stepper motors, it requires an intelligent controller (IC) to control them (Kuei-You, Kuan-Yi \& Chao-Chieh 2017). The steppers motors only require the number of steps and the directions to be programmed with an Intelligent Controller (IC) (Kuei-You, Kuan-Yi \& Chao-Chieh 2017). However, stepper motors start to slow down before it moves rapidly so it does not have any slippage for consistent moving (Kuei-You, Kuan-Yi \& Chao-Chieh 2017). For a robot, slippage is an important factor that needs to be considered. If the motor ever misses a step, the whole pseudo-code after would be disrupted due to one missed step which will lead to the robot failing to complete the whole track. Next, the materials to manufacture the wheels have to be considered. The wheels need to withstand a high amount of torque from the robot to move. Poly Lactic Acid (PLA) have a high Young’s Modulus and yield strength which is able to counter the high amount of torque that is being put from the weight of the robot (Farah, Anderson \& Langer 2016). Lastly, stability is an issue that the robot may face if it is driven with less number of wheels. Hence, omni-directional wheels are considered to provide support to the robot and to the rotation of the wheels. Omni-directional wheel consists of many small discs in a wheel (Islamgozhayev et al. 2015). The wheels can move or slide in any directions due to its disc hence it can move immediately in arbitrary directions (Islamgozhayev et al. 2015).


As a result of researching, it shows that each component plays different roles but they all coordinates with one and another to allow movement. The constraints and specifications also have to be considered in manufacturing the wheels.


%----------------------------------------------------------------------------------------
%Alternative designs [WIP]
%----------------------------------------------------------------------------------------
\chapter{Alternative designs}
\section{Design A - 8 wheels design}
After brainstorming through the initial stages of the robot building design procedure, the team had initially planned on making the robot with 8 wheels to achieve the aims of the competition in a timely manner {2 sets of 4 wheels in which one of the sets operate forwards/backwards while the motion of the other set of wheels operate left/right}. The robot mainly consists of a scissor lift, an extendable incline, a box capable of storing the power-pack and a base mounted on 8 wheels with 4 axles (each axle supporting 2 wheels) to hold the structure and electronics comprised in programming the robot. 
However, the 8–wheels design has some advantages and disadvantages which attribute to the robot’s core functionality.  The advantages includes high structural integrity to support the infrastructure on top of the base of the robot, total time accumulated to finish the course decreases, less programming is required to move the robot which was one of the main reasons that this design was proposed initially; the robot needs to successfully move the power-pack with the balls in place into the bunker along with the arrow on the power-pack pointing the same direction as it was picked up into the bunker for maximum points. For this, the team thought rotating the robot in a desired way will consume a lot of time as well as money to buy several motors and motor shields. Thus, the idea to build an 8–wheels design where no rotations is needed, and the robot primarily moves with 2 degrees of freedom (horizontal and vertical directions [x-y axis]) which does the job efficiently and effectively such that the power pack’s direction does not change during any part of the process. The power pack is collected via the extended arms thus probing the robot to change transmission allowing it to move towards the bunker. Once close to the bunker the original transmission wheels will move to close the distance and the motor on the of the box will activate rolling the incline plane down releasing an extended length of surface allowing the power pack to slide down precisely into the bunker with the arrow in the right direction. 
After more planning and advice from experienced people in the field of robotics, it was discovered that to do this, a lot of money would be required in order to build a transmission system such that the robot switches between the 2 sets of wheels during the course of its journey which also requires a lot of physical work, hence, the team decided that this design shall be disregarded due to the nature of difficulty in building this system along with its incorporated costs.

%%Insert Pictures

%===================Design B=============================================
\section{Design B - 4 Wheel design}
This robot is a system with a free moving arm which acts as both the lifting and the grabbing
mechanism. The robot mainly consists of an extended arm comprised of two sections so that it is
long enough to maneuver around the pole and collect the power pack and to also place it into the
bunker, a fork shaped platform to latch onto the payload whilst keeping the power pack in place
and a base mounted on 4 wheels with 2 stepper motors and two omni-directional wheels to hold
the structure and electronics comprised in programming the robot.
The 4-wheel arm system is more mobile and if executed will enable us more mobility to
approach the power pack. It also enables easier placement of power pack into the bunker as it is
able to rotate over the pole and position itself into bunker with precision. It is also able to rotate
thus giving it more options to go in several directions. The 4-wheel arm system is also light-
weight as it requires less raw material. However, this system will require more motors and motor
shields to operate the individual parts thus increasing the cost of the robot. The increased
complexity in programming is also a disadvantage as several rotations, movements and
maneuvers (with the arm) will need to be executed with high precision especially without the aid
of sensors with the current design. The electronics will also be an issue as connecting each parts
without them interfering with each other will need to be considered so that maximum efficiency
will be achieved.
The hands are extended and will rotate over the pole and sweeping the power pack into the
shaped platform, and will rotate over to the bunker side so that it can be placed with the arrow in
its initial position. Overall, the 4-wheel system provides good structural integrity, high mobility

and good maneuverability and hence is a more viable design which can be implemented in order
to achieve the aims of the competition.
%%Insert Pictures

%===================Design C=============================================
\section{Design C - 2 Wheels Design}
The 2-wheel robot system is a combination of the 8-wheels robot and the 4-wheel robot. Like the 8-wheel robot it has an extended claw mechanism to pick up the power pack and store it in a safe and secure space while waiting to be deposited into the bunker. Like the 4-wheels design however it has great maneuverability as it uses two stepper motors allowing it to rotate and move in any direction required.

 However, the 2 –wheel robot system has some advantages and disadvantages which attribute to the robot’s core functionality.  The disadvantages include less structural integrity to support the infrastructure on top of the base of the robot as it is not balanced and will tip over on one side. This will cause uncertainty in being able to keep the golf balls in the power pack as it is unstable. The total time accumulated to finish the course increases as well due to increased movement, rotations and also requires more torque to move as there will be drag from one side of the base touching the ground. It will also require some complex programming due to the movements and rotations that is required to complete the journey. Advantages of this system include less weight and less cost as only 2 stepper motors will be required and less maintenance issues as opposed to the other design proposals. It will also be easier to manufacture as it is a simple design that is able to be built using everyday materials.
The 2-wheels robot will collect the payload allowing it to slide at an angle and be blocked by a wall attached to a wire. It will rotate and travel towards the bunker. Once it reaches the bunker the robot will rotate to orientate itself so that the wall is facing the bunker and the motor will turn, lowering the wall allowing the power pack to slide into the payload.
Hence, considering the pros and cons of the design, it was disregarded due to its low reliability in stability as keeping the payload in place is crucial and worries of damage to the stepper motors while dragging the base through the floor is a point to consider as well.

%%Insert picture [WIP]
%%Insert table [WIP]

%----------------------------------------------------------------------------------------
%Final Design[WIP]
%----------------------------------------------------------------------------------------
\chapter{Final design}
\section{Design Choice}

There are many factors to consider while choosing the most feasible design to build. They include cost, stability, maneuverability, power draw, speed and lead time.


%Lead time = time for manufacturing... just more fancy industrial words
\begin{table}[htbp]
\centering
\begin{tabular}{|l|l|l|l|l|}\hline
Factors         & Weighting (\%) & 8 Wheel Design & 4 Wheel Design & 2 Wheel Design \\ \hline
Speed           & 10             & 8              & 6              & 6              \\ \hline
Cost            & 25             & 4              & 5              & 8              \\ \hline
Lead Time       & 30             & 3              & 7              & 8              \\ \hline
Stability       & 15             & 9              & 10             & 6              \\ \hline
Maneuverability & 15             & 9              & 5              & 7              \\ \hline
Power Draw      & 5              & 5              & 6              & 8              \\ \hline
Total           & 100            & 5.5            & 6.5            & 7.4           \\ \hline
\end{tabular}
\caption{Comparison Matrix}
\end{table}
%PLEASE DONT TOUCH TABLE... ANYTHING WRONG WITH IT ASK PH TO CHANGE.

A more detailed calculation of the total score, explanations for the factors and scores chosen are provided in the appendices (Refer to Appendix B). From the matrix, it is a clear choice to choose the 2 wheel design.


\begin{center}
\includegraphics[width = 0.9 \textwidth]{cad1.png}\\
Figure 1 : CAD Model of the robot
\end{center}


The final design of the robot is shown in Figure 1 which consists of a base with 2 wheels controlled by 2 stepper motors and 2 omnidirectional wheels to provide structural integrity for the robot. The base of the device is made out of cardboards stacked together as it is cheap and sufficiently tough for the robot. The wheels are 3D printed using Poly Lactic Acid (PLA) as it has similar properties to Acrylonitrile Butadiene (ABS) and is able to withstand a considerable amount of torque applied by the motor. The other aspect of the design of the robot contains a scissor lift which is made out of galvanized steel and controlled by a linear actuator. The scissor lift enables the system to provide a lifting mechanism such that it could reach the desired height of 800mm in order to collect the power-pack. In addition, there is also a collecting mechanism (mounted on top of the lift) consisting two arms along with a slide for an allocated storage space while the robot is in motion which is also connected to the top of the scissor lift. The collecting mechanism containing the two metal arms are operated by 2 DC motors in order to collect the power pack from Pole C. Moreover, the two metals arms are needed to hold the power-pack in place before delivering it into the bunker.

%Insert figure

\section{Machining}
The robot’s base is made of MDF standard wood, this was by far the cheapest option and also of less weight, however; its not great in durability, thus, the base was bound to another MDF sheet to provide structural integrity. The wheels were made of plastic and were 3D printed which were attached to the shaft of the stepper motors (attached underneath the base using brackets) which is held tight by a grub screw such that the wheels are locked in place at all times. The scissor lift, and the slide were made of galvanised steel metal sheets. Two brackets are mounted on the base while the scissor lift is attached to the L- shaped brackets with one end of the lift left free-to-slide such that the lift is operable upwards and downwards using the linear actuator. The pop sticks on the scissor lift were handcrafted by using workshop machinery including guillotine, metal benders and drill presses. The popsticks were double-folded to provide more strength and rigidity to the lift. On the other hand; tools such as tin-snippers, battery operated drills, files and hacksaws were used in making the lift as well as the slide. The popsticks were connected using tapered connecting metal rods which are bolted on both ends.The flaps controlled by the DC motors mounted on top of the lift is made of steel blocks with small tapered holes for grub screw fittings such that the flaps are locked to the shaft of the motors at all times. 

\section{3D printing}
The wheels were 3D printed with PLA as both ABS and PLA holds the same properties. The diagram in Figure 3 was the initial idea of the wheels. It has an axle which can be connected to a cover shaft which is also 3-D printing, holding both the axle of the wheel and the shaft of the stepper motors in place. However, it was not a solid design and was proven that the axle would snap with little amount of force applied to it. Hence, the design had to be reiterated as shown in Figure 4.
%insert figure

After a few consultations, the connection between the wheels and the shaft of the motor revolves around an extrusion and grub screw method. The idea of the wheel-motor assembly was to make an extrusion out of the wheel having a tight grip on the shaft of the motor. As shown in Figure 5, the motor shaft is placed in the hole of the extrusion and locked by an M-8 grub screw. 

\section{Electronics}
%Insert table
There are a few relays/motor shields controlling each motor. There were 2 big easy drivers, 1 4-channel relay module and 1 motor shield connected to the Arduino. The 2 Big Easy Drivers are connected to the 2 steppers motors respectively. Red and yellow wires from the steppers are connected to Motor A for the Big easy driver. Green and grey wires from the steppers are connected to Motor B for the other Big easy driver. Both sets of wires are paired up due to the similar resistances between the wires which are about 35Ω.  The 4-channel relay module is connected to the linear actuator which controls the scissor lift. Lastly, the motor shield connects to 2 DC motors which are mounted on top of the base controlling the two metal arms. All Vin on all the motor shields are connected to the positive supply rail and all of the ground of the motor shield is connected to the negative supply rail on a mini breadboard.

\section{Coding}
The coding is done in source code editor Atom and compiled in Arduino compiler. The reason behind having different editor and compiler is due to personal preference and Atom having more features such as brace matching, auto completion and code folding.

On the code itself, it uses a lot of accelstepper library functions. Arduino's standard stepper.h library is not used as the accelstepper contains more features including multiple stepper control, stepper acceleration control and stepper direction control. Those features are needed to simplify the code of the robot. 

While the code is in C++, the language's Object Oriented Properties are not really used to its full extent. The only objects called in the code is rightMotor and leftMotor, the bare minimum needed for accelstepper to function properly. While Object Oriented Programming does produce a more robust code, the robot does not require such robust programming as it is completely nonreactive to the outside environment. The robot does not have a sensor suite, and uses dead reckoning to navigate.

\section{Testing}
Testing phase was delayed until the final day before the competition due to several factors. Firstly, the wheels took a huge amount of time to 3D print. Approximately 6 hours were used to print each wheel and the design of the wheels had to be reiterated several times. In addition, there were nozzle jams while printing the wheels. Next, the manufacturing phase needing both stepper motors and DC motors had decreased the amount of time for testing phase. The robot had to be completely manufactured in order to test the motors. This in turn had delayed testing time of the robot. Finally, there was no allocated member in doing electronics. Due to this procrastination, it had also delayed testing time. The electronics were delayed to the final week. Even with the delay, the circuit had to be reconfigured at least 9 times as the circuit could not be carried around easily and due to some misunderstanding on the connection of the drivers. There was some uncertainty on the codes as well, as the code compiles however when it is uploaded there was no movement from any motors. This may also be due to incorrect connections in the circuits.

\section{Modifications}
There were a few modifications that needed to be done. First and foremost, some tolerances on the robot had to be changed. For example, the scissor lift operates when the linear actuator extends moving the connecting rod connected to the scissor lift. Due to certain miscalculated tolerances, the connecting rod could not slide through the slot which preventing the scissor lift from extending or retracting. Next, electronics instead of having to buy drivers or motor shields which are general and expensive. There were much simpler one-function motor shields such as stepper modules which are much cheaper in comparison. Furthermore, there were a few loose connections after reconfiguring the circuits a few times. Hence, the connections of the wires should be considered while doing the electronics as it may be a factor of the robot not moving. 
 
\pagebreak

%----------------------------------------------------------------------------------------
%Conclusion
%----------------------------------------------------------------------------------------
\chapter{Conclusion}
Three different designs for the movement of the robot and a scissor lift design was considered for this robot. A hybrid design of 2 driver wheel and 2 idle wheel were chosen as they offer the best trade off between maneuverability and stability. A scissor lift design was chosen as it is a relatively simple design to manufacture.

While the robot failed to achieve its objectives on the competition date, it is a proven concept that should work given enough time for field testing. However, it is undeniable that more optimization like increasing the voltage supplied to the linear actuator to reducing the lift extension period can be done on the robot. By doing so, a better and more efficient robot can be built.
\pagebreak

%----------------------------------------------------------------------------------------
%Reference List
%----------------------------------------------------------------------------------------


\chapter{Reference List}
\begin{flushleft}
Farah S., Anderson D. G. \& Langer R. 2016, ‘Physical and mechanical properties of PLA, and their functions in widespread applications — A comprehensive review’, Advanced Drug Delivery Reviews, Volume 107.
\linebreak
\linebreak
Greenough N. \& Kung C. C. 2013, ‘A new high-efficiency stepper motor driver for old technology stepper motors’, in 2013 IEEE 25th Symposium on Fusion Engineering (SOFE), San Francisco, CA.
\linebreak
\linebreak
Islamgozhayev, T. U., Mazhitov, S.S., Zholmyrzayev, A.K. and Toishybek, E.T. 2015, ‘IICT-bot: Educational robotic platform using omnidirectional wheels with open source code and architecture’, in 2015 International Siberian Conference on Control and Communications (SIBCON).
\linebreak
\linebreak
Kuei-You L, Kuan-Yi W \& Chao-Chieh L 2017, ‘High-performance series elastic stepper motors for interaction force control’, in 2017 IEEE International Conference on Advanced Intelligent Mechatronics (AIM), Munich.
\linebreak
\linebreak
Warren, J.D., Adams, J. \& Molle, H. 2011, ‘Arduino Robotics’, Collection: Technology in Action. New York (USA): Springer Verlag GmbH.
\end{flushleft}
\pagebreak


\section*{Appendix A}


\end{document}
